% ================================================================
% CHAPTER 6: PAYOFFS AS INNER PRODUCTS
% ================================================================
\chapter{Payoffs as Inner Products}
\label{ch:payoffs}

\epigraph{In the beginning was the Word, and the Word was with God, and the
Word was God.}{John 1:1}

\section{The Meaning Game}

We now make a fundamental shift in perspective. In Part~I, ideas were objects
of contemplation: we studied their geometry, their distances, their resonances.
In Part~II, ideas become \emph{strategic instruments}: they are signals sent by
agents who seek to influence the beliefs and actions of others. The Hilbert
space structure of Part~I provides the geometry; game theory provides the
strategic logic.

The basic setup is simple. There are two agents --- a \textbf{sender} and a
\textbf{receiver} --- each holding an idea. The sender transmits a signal, and
the receiver interprets the signal by composing it with their own prior belief.
The payoff to each agent is the resonance between their own idea and the
resulting interpretation.

\begin{definition}[The Meaning Game]
\label{def:meaning-game}
A \textbf{meaning game} consists of:
\begin{enumerate}
\item A Hilbert ideatic space $(\IS, \compose, \void, \Hspace, \varphi)$;
\item A \textbf{sender} with idea $s \in \IS$ (the \emph{signal});
\item A \textbf{receiver} with idea $r \in \IS$ (the \emph{prior belief});
\item The \textbf{interpretation}: the receiver's updated belief after receiving
the signal is $r' = r \compose s$ (the composition of their prior with the
signal);
\item The \textbf{sender payoff}: $\payoff{S}{s, r} := \rs{s}{r \compose s}$;
\item The \textbf{receiver payoff}: $\payoff{R}{s, r} := \rs{r}{r \compose s}$.
\end{enumerate}
\end{definition}

The payoff structure is motivated by the idea that each agent benefits from
the interpretation (the updated belief $r \compose s$) to the extent that it
resonates with their own idea. The sender wants the interpretation to resonate
with the signal they sent; the receiver wants the interpretation to resonate
with their prior belief.

\begin{remark}[Why These Payoffs?]
\label{rmk:payoff-motivation}
The choice of payoff $\rs{s}{r \compose s}$ rather than, say,
$\rs{s}{r}$ (direct resonance between signal and prior) deserves comment.
The direct resonance $\rs{s}{r}$ measures pre-existing agreement between
sender and receiver, which is not under either agent's control. The payoff
$\rs{s}{r \compose s}$ measures how well the \emph{interpreted} message
resonates with the sender's intention --- and this \emph{does} depend on the
signal choice. It is the resonance between what was meant and what was
understood.
\end{remark}


\section{The Payoff Decomposition}

The key result of this chapter decomposes each agent's payoff into two terms
with clear interpretations.

\begin{theorem}[Sender Payoff Decomposition]
\label{thm:sender-payoff}
The sender's payoff decomposes as:
\[
  \payoff{S}{s, r} = \rs{s}{s} + \rs{s}{r}.
\]
That is: sender payoff $=$ self-resonance $+$ cross-resonance.
\end{theorem}

\begin{proof}
\begin{align*}
  \payoff{S}{s, r}
  &= \rs{s}{r \compose s} \\
  &= \ip{\embed{s}}{\embed{r \compose s}} \\
  &= \ip{\embed{s}}{\embed{r} + \embed{s}} \\
  &= \ip{\embed{s}}{\embed{r}} + \ip{\embed{s}}{\embed{s}} \\
  &= \rs{s}{r} + \rs{s}{s}.
\end{align*}

In Lean~4:
\begin{lstlisting}
theorem sender_payoff (s r : Idea) :
    rs s (r ∘ s) = rs s s + rs s r := by
  simp [rs, embed_compose, inner_add_right]
\end{lstlisting}
\end{proof}

\begin{theorem}[Receiver Payoff Decomposition]
\label{thm:receiver-payoff}
The receiver's payoff decomposes as:
\[
  \payoff{R}{s, r} = \rs{r}{r} + \rs{r}{s}.
\]
That is: receiver payoff $=$ self-resonance $+$ cross-resonance.
\end{theorem}

\begin{proof}
\begin{align*}
  \payoff{R}{s, r}
  &= \rs{r}{r \compose s} \\
  &= \ip{\embed{r}}{\embed{r} + \embed{s}} \\
  &= \ip{\embed{r}}{\embed{r}} + \ip{\embed{r}}{\embed{s}} \\
  &= \rs{r}{r} + \rs{r}{s}.
\end{align*}
\end{proof}

\begin{interpretation}[The Anatomy of Communicative Success]
\label{interp:payoff}
The payoff decomposition is one of the most striking results in IDT~v2. It
reveals that each agent's payoff from communication has \emph{exactly two
components}:

\begin{enumerate}
\item \textbf{Self-resonance} ($\rs{s}{s}$ or $\rs{r}{r}$): your own conviction
--- the weight you bring to the interaction. This is independent of the other
agent; it is your ``baseline'' payoff from holding your idea.

\item \textbf{Cross-resonance} ($\rs{s}{r}$ or $\rs{r}{s}$): the degree of
pre-existing agreement between sender and receiver. This is symmetric by
Theorem~\ref{thm:resonance-symmetry}: the cross-resonance benefit is the
\emph{same} for both agents.
\end{enumerate}

The immediate consequence is a \textbf{fundamental insight about communication}:
\emph{your payoff from communication equals your own conviction plus how much
the other person already agrees with you.} Communication benefits both parties
equally from the ``agreement'' component ($\rs{s}{r}$ appears in both payoffs),
and additionally rewards each party in proportion to their own conviction.

This has profound implications for the dynamics of discourse:

\paragraph{Echo chambers.} In an echo chamber, $\rs{s}{r} \gg 0$: sender and
receiver already agree strongly. Both get large payoffs, reinforcing the
behavior of communicating within the echo chamber. The payoff formula explains
why echo chambers are so psychologically rewarding: not only do you feel good
about your own ideas ($\rs{s}{s}$), but the other person's agreement provides
an \emph{additional} boost ($\rs{s}{r}$).

\paragraph{Adversarial communication.} When $\rs{s}{r} < 0$, the
cross-resonance term is negative. The sender's payoff can still be positive
(if $\rs{s}{s} > |\rs{s}{r}|$), but it is reduced by the receiver's
disagreement. This explains why arguing with hostile audiences is costly even
when your own conviction is high: the negative cross-resonance erodes your
payoff.

\paragraph{Preaching to the void.} When $r = \void$ (the receiver has no prior
belief), $\rs{s}{\void} = 0$ and the sender's payoff reduces to $\rs{s}{s}$
--- pure self-resonance. You benefit only from your own conviction; there is
no agreement to amplify it. This is the mathematical model of ``talking to
oneself'' or ``shouting into the void.''
\end{interpretation}

\begin{theorem}[Total Payoff]
\label{thm:total-payoff}
The \textbf{total payoff} of the meaning game is:
\[
  \payoff{S}{} + \payoff{R}{} = \rs{s}{s} + \rs{r}{r} + 2\rs{s}{r}
  = \rs{s \compose r}{s \compose r} = w(s \compose r)^2.
\]
\end{theorem}

\begin{proof}
Adding the two payoff decompositions:
\begin{align*}
  \payoff{S}{} + \payoff{R}{}
  &= (\rs{s}{s} + \rs{s}{r}) + (\rs{r}{r} + \rs{r}{s}) \\
  &= \rs{s}{s} + \rs{r}{r} + 2\rs{s}{r} \\
  &= w(s \compose r)^2,
\end{align*}
using the norm expansion formula (Theorem~\ref{thm:norm-expansion}).

In Lean~4:
\begin{lstlisting}
theorem total_payoff (s r : Idea) :
    rs s (r ∘ s) + rs r (r ∘ s) = rs (s ∘ r) (s ∘ r) := by
  simp [rs, embed_compose, inner_add_left, inner_add_right]
  ring
\end{lstlisting}
\end{proof}

\begin{interpretation}[Communication as Norm Squared]
\label{interp:total}
The total payoff of the meaning game equals the \emph{squared norm} of the
composite idea $s \compose r$. This is a remarkable result: it says that the
total value of a communicative exchange equals the weight of the shared
understanding that results from it.

When the shared understanding is weighty ($w(s \compose r)$ is large), both
parties benefit. When it is light ($w(s \compose r)$ is small), neither benefits
much. And when the composition cancels ($s$ and $r$ are perfectly opposed,
yielding $w(s \compose r) = 0$), the total payoff is zero: the exchange
produces no value for either party.

This connects to Habermas's theory of \emph{communicative action}: successful
communication creates a shared ``lifeworld'' (Lebenswelt) of mutual
understanding. In IDT~v2, this lifeworld is the composite idea $s \compose r$,
and its weight $w(s \compose r)^2$ is the total value of the communicative
exchange. Habermas's normative ideal of ``undistorted communication'' would
correspond, in this model, to a game where both parties' signals are authentic
(they signal their true ideas) and the resulting composite has maximum weight.
\end{interpretation}


\section{Payoff Comparison: Sender vs.\ Receiver}

Who benefits more from communication --- the sender or the receiver?

\begin{proposition}[Payoff Difference]
\label{prop:payoff-diff}
\[
  \payoff{S}{} - \payoff{R}{} = \rs{s}{s} - \rs{r}{r} = w(s)^2 - w(r)^2.
\]
The sender benefits more if and only if $w(s) > w(r)$: the agent with the
weightier idea gets the higher payoff.
\end{proposition}

\begin{proof}
$\payoff{S}{} - \payoff{R}{} = (\rs{s}{s} + \rs{s}{r}) - (\rs{r}{r} + \rs{r}{s})
= \rs{s}{s} - \rs{r}{r}$, using $\rs{s}{r} = \rs{r}{s}$.
\end{proof}

\begin{interpretation}[The Advantage of Conviction]
\label{interp:conviction-advantage}
Proposition~\ref{prop:payoff-diff} reveals a stark truth: \emph{in the meaning
game, conviction is power}. The agent with the weightier idea (higher
self-resonance) always gets the higher payoff, regardless of the degree of
pre-existing agreement. The cross-resonance term $\rs{s}{r}$ is symmetric
and benefits both agents equally; only the self-resonance term differs, and it
favors the more convicted agent.

This formalizes a phenomenon well-known in rhetoric and politics: confident,
forceful communication (high $w(s)$) is rewarded disproportionately, independent
of its truth or wisdom. A demagogue with enormous conviction ($w(s) \gg w(r)$)
always ``wins'' the meaning game against a moderate with less conviction, even if
the moderate's ideas have more merit in some external sense. The Hilbert space
model captures the structural advantage of conviction without endorsing it
normatively.

In Bourdieu's terminology, self-resonance is a form of \emph{symbolic capital}:
the capacity to make one's ideas count in the cultural field. Agents with high
symbolic capital (high norm) benefit more from every interaction, creating a
self-reinforcing dynamic: those who start with more symbolic capital
accumulate more payoff, which (in a dynamic model) may translate into even
more capital. The ``rich get richer'' dynamic in cultural production has a
geometric foundation in the Hilbert space model.
\end{interpretation}


\section{Strategic Signal Choice}

We now consider the sender's strategic problem: given the receiver's prior $r$,
what signal $s$ maximizes the sender's payoff?

\begin{theorem}[Optimal Signal]
\label{thm:optimal-signal}
If the sender can choose any signal $s$ with $\nrm{\embed{s}} \leq M$ (a
``budget constraint'' on signal weight), the optimal signal is:
\[
  \embed{s^*} = M \cdot \frac{\embed{r}}{\nrm{\embed{r}}} + M \cdot \hat{e},
\]
where $\hat{e}$ is any unit vector. In particular, any signal pointing in the
direction of $\embed{r}$ achieves at least as much payoff as a signal of the same
weight pointing in any other direction.

More precisely, for fixed $\nrm{\embed{s}} = M$, the sender's payoff
$\payoff{S}{} = M^2 + M \cdot \nrm{\embed{r}} \cdot \cos\alpha$, where
$\alpha$ is the angle between $\embed{s}$ and $\embed{r}$. This is maximized
when $\alpha = 0$ (signal parallel to receiver's prior):
\[
  \payoff{S}{\max} = M^2 + M \cdot \nrm{\embed{r}}.
\]
\end{theorem}

\begin{proof}
The sender's payoff is $\payoff{S}{} = \rs{s}{s} + \rs{s}{r} = M^2 +
\ip{\embed{s}}{\embed{r}}$. For fixed $\nrm{\embed{s}} = M$, the term
$\ip{\embed{s}}{\embed{r}}$ is maximized when $\embed{s}$ is parallel to
$\embed{r}$ (by Cauchy--Schwarz), yielding $\ip{\embed{s}}{\embed{r}} = M \cdot
\nrm{\embed{r}}$. The maximum payoff is thus $M^2 + M \cdot \nrm{\embed{r}}$.

In Lean~4:
\begin{lstlisting}
theorem optimal_signal_payoff (hM : ‖φ s‖ = M)
    (hr : φ r ≠ 0) :
    rs s s + rs s r ≤ M ^ 2 + M * ‖φ r‖ := by
  simp [rs]
  have h1 : inner (φ s) (φ s) = M ^ 2 := by rw [real_inner_self_eq_norm_sq, hM]
  have h2 : inner (φ s) (φ r) ≤ M * ‖φ r‖ := by
    calc inner (φ s) (φ r) ≤ ‖φ s‖ * ‖φ r‖ := real_inner_le_norm _ _
    _ = M * ‖φ r‖ := by rw [hM]
  linarith
\end{lstlisting}
\end{proof}

\begin{interpretation}[The Panderer's Dilemma]
\label{interp:panderer}
The optimal signal theorem reveals a troubling strategic logic:
\emph{the optimal signal is always in the direction of the receiver's prior}.
To maximize payoff, the sender should tell the receiver what they want to hear.
This is the ``panderer's dilemma'': strategic communication incentivizes
flattering the receiver's beliefs rather than communicating truth.

The result is driven by the cross-resonance term $\rs{s}{r}$, which is maximized
when $\embed{s}$ is parallel to $\embed{r}$. A sender who sends a signal
orthogonal to the receiver's prior gets zero cross-resonance (no agreement
bonus); a sender who sends a signal opposed to the receiver gets negative
cross-resonance (an agreement penalty).

In political terms, this is a model of \emph{populism}: the optimal political
message (in terms of payoff) is the one that most closely mirrors the electorate's
existing beliefs, regardless of its truth or wisdom. The model does not predict
that all communication is pandering --- a sender may have reasons beyond payoff
to send a truthful signal --- but it shows that the \emph{strategic incentive}
always favors pandering.

The only escape from the panderer's dilemma, within this model, is to change the
payoff function. If the sender's payoff depended on some external criterion of
truth rather than resonance with the receiver, the incentive to pander would
vanish. This connects to the philosophical distinction between \emph{rhetoric}
(the art of persuasion, which rewards pandering) and \emph{dialectic} (the art
of truth-seeking, which penalizes it).
\end{interpretation}


\section{The Symmetric Meaning Game}

\begin{definition}[Symmetric Meaning Game]
\label{def:symmetric-game}
A meaning game is \textbf{symmetric} if both agents simultaneously choose
signals and interpret the other's signal. Agent $i$ has idea $a_i$ and receives
signal $a_j$ from agent $j$. Agent $i$'s payoff is:
\[
  \payoff{i}{} = \rs{a_i}{a_i \compose a_j} = \rs{a_i}{a_i} + \rs{a_i}{a_j}.
\]
\end{definition}

\begin{theorem}[Symmetric Game Is Not Zero-Sum]
\label{thm:not-zero-sum}
The symmetric meaning game is not zero-sum in general. The total payoff is:
\[
  \payoff{1}{} + \payoff{2}{} = \rs{a_1}{a_1} + \rs{a_2}{a_2} + 2\rs{a_1}{a_2}
  = w(a_1 \compose a_2)^2.
\]
This is positive unless $a_1 \compose a_2 \equiv_\varphi \void$.
\end{theorem}

\begin{proof}
Sum the payoffs: $(\rs{a_1}{a_1} + \rs{a_1}{a_2}) + (\rs{a_2}{a_2} +
\rs{a_2}{a_1}) = \rs{a_1}{a_1} + \rs{a_2}{a_2} + 2\rs{a_1}{a_2} =
w(a_1 \compose a_2)^2 \geq 0$.
\end{proof}

\begin{interpretation}[Communication Creates Value]
\label{interp:value-creation}
The non-zero-sum nature of the meaning game is fundamental. Communication is
not a zero-sum competition where one party's gain is the other's loss.
Rather, communication \emph{creates value}: the total payoff $w(a_1 \compose
a_2)^2$ is positive whenever the composite idea has non-zero weight. Both
parties benefit from the exchange (though not necessarily equally --- the
more convicted party benefits more, as shown in Proposition~\ref{prop:payoff-diff}).

The total value created is maximized when the ideas are perfectly aligned
($w(a_1 \compose a_2)^2 = (w(a_1) + w(a_2))^2$) and minimized when they are
perfectly opposed ($w(a_1 \compose a_2)^2 \approx 0$). Orthogonal ideas create
an intermediate amount of value ($w(a_1 \compose a_2)^2 = w(a_1)^2 + w(a_2)^2$
--- the Pythagorean theorem).

This connects to the economic theory of \emph{gains from trade}: just as
voluntary exchange creates surplus (both parties end up better off), voluntary
communication creates meaning (both parties end up with higher-weight
interpretations). The ``gains from communication'' are measured by
$w(a_1 \compose a_2)^2$ and are positive whenever the ideas do not completely
cancel.
\end{interpretation}

\begin{example}[Payoffs in Academic Discourse]
\label{ex:academic}
Consider an academic conference where three scholars present ideas embedded in
$\mathbb{R}^3$:
\begin{align*}
  \embed{a_1} &= (4, 0, 0) \quad \text{(specialist in dimension 1)}, \\
  \embed{a_2} &= (3, 3, 0) \quad \text{(bridging dimensions 1--2)}, \\
  \embed{a_3} &= (0, 0, 5) \quad \text{(specialist in dimension 3)}.
\end{align*}
Payoff matrix (sender payoff in each cell):
\begin{center}
\begin{tabular}{c|ccc}
$\payoff{S}{}$ & $r = a_1$ & $r = a_2$ & $r = a_3$ \\
\hline
$s = a_1$ & -- & $16 + 12 = 28$ & $16 + 0 = 16$ \\
$s = a_2$ & $18 + 12 = 30$ & -- & $18 + 0 = 18$ \\
$s = a_3$ & $25 + 0 = 25$ & $25 + 0 = 25$ & -- \\
\end{tabular}
\end{center}
Scholar $a_2$ (the bridge-builder) gets the highest payoff when addressing
$a_1$ (payoff 30), because they share dimension 1. Scholar $a_3$ gets equal
payoffs regardless of audience (payoff 25), because they are orthogonal to
everyone --- a mathematical model of the isolated specialist who can talk
to anyone equally (in)effectively.
\end{example}

\bigskip

This chapter has established the game-theoretic framework for communication in
Hilbert space: payoffs decompose into self-resonance (conviction) plus
cross-resonance (agreement), communication creates value, and the optimal
signal --- for purely strategic agents --- mirrors the receiver's prior. In
Chapter~\ref{ch:nash}, we study the equilibrium structure of these games.
